%------------------------------------------------------------------------------%
\section{Círculo Trigonométrico Unitário}

Considere o triângulo retângulo de hipotenusa 1,
apresentado na Figura~\ref{fig:trigon1}.

%------------------------------------------------------------------------------%
\begin{figure}
  \centering
  \includegraphics[width=0.6\textwidth,trim={3cm 13cm 3cm 15cm},clip]{figcirctrig1}
  \caption{Triângulo retângulo de hipotenusa 1.}
  \label{fig:trigon1}
\end{figure}
%------------------------------------------------------------------------------%

Suponhamos que
\begin{enumerate}
  \item o triângulo esteja situado no plano cartesiano,

  \item o vértice associado ao ângulo $\theta$ esteja sobre a origem $(0, 0)$,

  \item o cateto adjacente ao ângulo $\theta$ fique sobre a parte positiva do
    eixo $x$,

  \item  a hipotenusa do triângulo retângulo meça 1 e o ângulo $\theta$ seja
    medido no sentido anti-horário a partir do eixo $x$.
\end{enumerate}
{\color{red}
Conforme observação anterior.
}

%------------------------------------------------------------------------------%
\begin{figure}
  \centering
  \includegraphics[width=0.6\textwidth,trim={2cm 11cm 3cm 160mm},clip]{figcirctrig2}
  \caption{?}
  \label{fig:?}
\end{figure}
%------------------------------------------------------------------------------%

{\color{red}
Podemos variar os valores dos ângulos que para dois valores $\theta_1$ e
$\theta_2$, teremos (???????????)
}

%------------------------------------------------------------------------------%
\begin{figure}
  \centering
  \includegraphics[width=0.6 \textwidth,trim={3cm 165mm 2cm 6cm} ,clip]{figcirctrig3}
  \caption{?}
  \label{fig:trigon2}
\end{figure}
%------------------------------------------------------------------------------%

Podemos agora definir o círculo trigonométrico unitário
(conforme Figura~\ref{fig:trigon3}),
que é o círculo do plano cartesiano de centro na origem $(0, 0)$ e raio 1.
O eixo $x$ é também chamado de
\emph{eixo dos cossenos}\index{Eixo!dos cossenos}, já o eixo
$y$ é também chamado de
\emph{eixo dos senos}\index{Eixo!dos senos}.

Considere um ponto $A$ sobre o círculo trigonométrico no primeiro quadrante e
seja $B$ a projeção desse ponto no eixo $x$.  Seja $C$ a interseção do
segmento de reta que passa por $\overline{OA}$ com a reta $x=1$ e $D=(1, 0)$, ou
seja, o segmento $\overline{CD}$ encontra-se sobre a reta $x=1$. Seja $\theta$
o ângulo $A\hat{O}B$. Observe que as coordenadas de $A$ no plano cartesiano são
$(\cos\theta, \sin\theta)$. Por semelhança de triângulos (triângulos $AOB$ e
$COD$) vemos que
\[
    \frac{\,\abs{\overline{CD}}\,}{1}
  = \frac{\,\abs{\overline{CD}}\,}{\abs{\overline{OD}}}
  = \frac{\,\abs{\overline{AB}}\,}{\abs{\overline{OB}}}
  = \frac{\,\sin\theta\,}{\cos\theta},
\]
ou seja,
\[
  \abs{\overline{CD}} = \tan\theta.
\]

%------------------------------------------------------------------------------%
\begin{figure}
  \centering
  \includegraphics[scale=3]{circtrigonom}
  \caption{Círculo Trigonométrico Unitário.}
  \label{fig:trigon3}
\end{figure}
%------------------------------------------------------------------------------%

Vemos assim que para $\theta\in(0, \pi/2)$ os valores $\sin\theta$, $\cos\theta$
e $\tan\theta$ podem ser obtidos a partir das coordenadas dos pontos $A$ e $C$:
$A=(\cos\theta, \sin\theta)$ e $C=(1, \tan\theta)$. Mas essa construção pode ser
feita para qualquer $\theta\in\R$ e dessa forma podemos definir as funções
seno, cosseno e tangente para quaisquer ângulos reais.  Por exemplo, essa
construção para um ângulo $\theta$ no segundo quadrante pode ser vista na
Figura~\ref{fig:trigon4}.

%------------------------------------------------------------------------------%
\begin{figure}
  \centering
  \includegraphics[width=0.4\textwidth,trim={0 1cm 9cm 3cm},clip]{circ5}
  \caption{Interpretação Geométrica da Tangente de um Ângulo no Segundo Quadrante}
  \label{fig:trigon4}
\end{figure}
%------------------------------------------------------------------------------%

Continuando a construção, podemos dividir o círculo trigonométrico unitário em
quatro partes iguais, chamadas de quadrantes. No primeiro quadrante estão os
ângulos entre 0 e \ang{90}, ou equivalentemente, entre 0 e
$\sfrac{\pi}{2}$. No segundo quadrante estão os ângulos entre \ang{90} e
\ang{180} (ou entre $\sfrac{\pi}{2}$ e $\pi$). No terceiro quadrante estão
os ângulos entre \ang{180} e \ang{270} (ou entre $\pi$ e
$\sfrac{3\pi}{2}$). Por fim, no quarto quadrante estão os ângulos entre
\ang{270} e $3\ang{60}$ (ou entre $\sfrac{3\pi}{2}$ e $2\pi$). Os ângulos
$0,\sfrac{\pi}{2}, \pi, \sfrac{3\pi}{2}$ e $2\pi$ são denominados ângulos do
eixo. Esses ângulos são notáveis e observando a Figura~\ref{fig:trigon3}, podemos
descobrir o valor do seno, cosseno e tangente desses arcos.
A Tabela a seguir apresenta esses valores.

%------------------------------------------------------------------------------%
\begin{ceqn}
\[
  \begin{array}{l<{\quad}>{\quad}c<{\quad}>{\quad}c<{\quad}>{\quad}c<{\quad}>{\quad}c<{\quad}>{\quad}r}
    \hline
            & 0 & \sfrac{\pi}{2} & \pi & \sfrac{3\pi}{2} & 2\pi \\ \hline
    \sin(x) & 0 & 1              & 0   & -1              & 0    \\
    \cos(x) & 1 & 0              & -1  & 0               & 1    \\
    \tan(x) & 0 & \nexists       & 0   & \nexists        & 0    \\ \hline
  \end{array}
\]
\end{ceqn}
%------------------------------------------------------------------------------%

%------------------------------------------------------------------------------%
\begin{observation}
  A \emph{relação fundamental trigonométrica}\index{Relação fundamental trigonométrica}
  A partir da Figura~\ref{fig:trigon3}, por exemplo, podemos extrair uma relação
  trigonométrica, que é muito útil. Essa relação é conhecida como
  relação fundamental trigonométrica e pode ser demonstrada, aplicando
  o Teorema de Pitágoras ao triângulo de hipotenusa 1 e catetos $\sin(\theta)$
  e $\cos(\theta)$
  \[
    \sin^2(x)+ \cos^2(x)=1
  \]
  A partir dessa relação podemos obter duas outras conforme serão mostradas nos
  exercícios que seguem.
\end{observation}
%------------------------------------------------------------------------------%

%------------------------------------------------------------------------------%
\begin{example}
  Usando a relação fundamental trigonométrica mostre que:
  \[
    \tan^2(x)+1 = \sec^2(x),
  \]
  quando
  \[
    \cos(x) \neq 0
    \quad\text{ e }\quad
    1+ \cot^2(x) = \csc^2(x)
    \quad\text{ se }\quad
    \sin(x) \neq 0.
  \]

  \solution
  Da relação fundamental trigonométrica temos:
  \begin{equation} \label{eq:relafun}
    \sin^2(x)+\cos^2(x)=1
  \end{equation}
  Sabendo que $\cos(x) \neq 0$, podemos dividir a equação~\eqref{eq:relafun}
  por $\cos^2(x)$. Assim,
  \[
    \frac{\sin^2(x)}{\cos^2(x)}+ 1= \frac{1}{\cos^2(x)}
  \]
  Logo,
  \[
    \tan^2(x)+1=\sec^2(x).
  \]
  Para chegarmos na outra expressão basta dividirmos a relação fundamental por
  $\sin^2(x)$.
\end{example}
%------------------------------------------------------------------------------%
